\documentclass{article}
\usepackage{PRIMEarxiv} % Core package for the PrimeAI Template
\usepackage{amsmath, amssymb, amsthm, bm, dcolumn} % Add amsthm here for the proof environment
\usepackage[numbers,sort&compress]{natbib} % Natbib for citations
\usepackage{graphicx} % For high-quality images
\usepackage{multicol} % Multi-column figures/tables
\usepackage{hyperref} % Hyperlinks
\usepackage{listings} % Code listings
\usepackage{authblk} % For structured affiliations
% Define theorem style (optional)
\theoremstyle{plain}
\newtheorem{theorem}{Theorem}
\renewcommand{\qedsymbol}{}

% Custom commands for primes and qubit encoding
\newcommand{\primeQubit}[1]{|p_{#1}\rangle}
\newcommand{\primeGate}[1]{U_{p_{#1}}}

\usepackage{wrapfig}
\usepackage[pscoord]{eso-pic}
\usepackage[fulladjust]{marginnote}
\reversemarginpar

% Typesetting improvements without footnote patching
\usepackage[protrusion=true, expansion=true, tracking=false]{microtype}
\microtypecontext{spacing=nonfrench}

% Line numbers
\usepackage[right]{lineno}

% Text layout - adjust as needed
\raggedright
\setlength{\parindent}{0.5cm}
\textwidth 5.25in 
\textheight 8.75in

% Set double spacing
\usepackage{setspace} 
\doublespacing

% Adjust width for specific content
\usepackage{changepage}

% Adjust caption style
\usepackage[aboveskip=1pt,labelfont=bf,labelsep=period,singlelinecheck=off]{caption}

% Remove brackets from references
\makeatletter
\renewcommand{\@biblabel}[1]{\quad#1.}
\makeatother

% Header, footer, and page numbers
\usepackage{lastpage,fancyhdr}
\pagestyle{fancy}
\fancyhf{}
\fancyhead[L]{Preprint - PrimeAI Enhanced Template}
\fancyfoot[C]{\scriptsize Multiplicity Theory © 2024 Ryan Van Gelder - Citizen Gardens \\ Licensed Under MIT and CC BY-NC-SA 4.0.}
\fancyfoot[R]{Page \thepage\ of \pageref{LastPage}}
\renewcommand{\footrule}{\hrule height 2pt \vspace{2mm}}

\title{\textbf{The Universal Multiplicity Constant ($\Lambda_m$)}: \\ A Foundational Mathematical Framework for Existence}
\author{Citizen Gardens - The Foundation of Multiplicity \\ \texttt{info@citizengardens.org}}
\date{March 06, 2025}

\begin{document}

\maketitle

\begin{abstract}
This paper introduces the Universal Multiplicity Constant (\(\Lambda_m\)) as an extension of Einstein’s General Relativity through Prime-Indexed Recursive Tensor Mathematics (PIRTM). Replacing the static cosmological constant (\(\Lambda\)) with a dynamic, prime-indexed tensor structure, we derive modified field equations encoding recursive spacetime dynamics. Enhanced by functorial mappings, homotopy fibrations, and SU(10) symmetry, \(\Lambda_m\) bridges quantum gravity, gravitational wave propagation, and black hole physics, offering testable predictions via LIGO, event horizon telescopes, and CMB analysis (Page 1).
\end{abstract}

\section{Fundamental Physics: Tensor Evolution of Spacetime}
Einstein’s field equations describe spacetime curvature:
\begin{equation}
R_{\mu\nu} - \frac{1}{2} g_{\mu\nu} R + \Lambda g_{\mu\nu} = 8\pi G T_{\mu\nu}.
\end{equation}
Their incompatibility with quantum mechanics prompts the introduction of \(\Lambda_m\), a recursive tensor correction within PIRTM (Section 1).

\subsection{The Extended Einstein Equations}
\(\Lambda_m\) is defined dynamically:
\begin{equation}
\Lambda_m = \sum_{p_i} \alpha_i p_i^{\beta_i} T_{ij}^{(p_i)},
\end{equation}
where:
\begin{itemize}
    \item \( p_i \in P = \{2, 3, 5, \dots\} \) are prime indices (Section 2.1).
    \item \( \alpha_i = \frac{1}{\log p_i} \) and \( \beta_i \) are scaling factors.
    \item \( T_{ij}^{(p_i)} \) are prime-indexed tensors, evolved recursively (Section 1.2).
\end{itemize}
The modified field equations are:
\begin{equation}
R_{\mu\nu} - \frac{1}{2} g_{\mu\nu} R + \Lambda_m g_{\mu\nu} = 8\pi G T_{\mu\nu}^{\text{recursive}},
\end{equation}
with:
\begin{equation}
T_{\mu\nu}^{\text{recursive}} = \sum_{k} \Lambda_m T_{ij}^{(p_k)},
\end{equation}
stabilized by spectral fixed points \( T^{(\infty)} = \frac{F}{1 - k} \) (Section 3).

\subsection{Effects on Gravitational Waves}
The standard wave equation \( \Box h_{\mu\nu} = 0 \) becomes:
\begin{equation}
\Box h_{\mu\nu} = \Lambda_m h_{\mu\nu},
\end{equation}
with plane wave solution \( h_{\mu\nu} = A_{\mu\nu} e^{i (k_\alpha x^\alpha)} \):
\begin{equation}
\eta^{\alpha\beta} k_\alpha k_\beta = \Lambda_m,
\end{equation}
yielding:
\begin{equation}
\omega^2 - |\vec{k}|^2 = \Lambda_m.
\end{equation}
\(\Lambda_m > 0\) suggests superluminal propagation, testable via LIGO (Page 49), with homotopy fibrations (Section 4.1) ensuring phase coherence.

\subsection{Effects on Black Holes}
The Schwarzschild metric adjusts to:
\begin{equation}
ds^2 = -\left(1 - \frac{2GM}{r} - \frac{\Lambda_m r^2}{3}\right) dt^2 + \left(1 - \frac{2GM}{r} - \frac{\Lambda_m r^2}{3}\right)^{-1} dr^2 + r^2 d\Omega^2,
\end{equation}
shifting the event horizon:
\begin{equation}
r_s = \frac{2GM}{c^2} + \mathcal{O}(\Lambda_m).
\end{equation}
Hawking temperature becomes:
\begin{equation}
T_H = \frac{\hbar c^3}{8\pi G M} \left(1 + \frac{\Lambda_m}{3}\right),
\end{equation}
with SU(10) symmetry enhancing stability (Page 13).

\subsection{Experimental Validation}
Proposed tests include:
\begin{itemize}
    \item Gravitational wave dispersion via LIGO.
    \item Black hole radius deviations via Event Horizon Telescope.
    \item CMB fluctuations from recursive spacetime, per PIRTM’s fractal evolution (Section 4.3).
\end{itemize}

\subsection{Conclusion}
\(\Lambda_m\) unifies quantum mechanics and gravity within PIRTM, offering testable deviations from classical relativity (Page 48).


\nocite{*}
\bibliographystyle{unsrt}
\bibliography{references}

\end{document}